\documentclass[11pt]{article}
\usepackage{explorer}
\newcommand{\ExplorerLanguage}{FR}
\newcommand{\ExplorerNumber}{07}
\newcommand{\ExplorerKind}{Exercices progressifs}
\begin{document}
\ExplorerTitle{Première--Terminale · probabilités}{Inverser une probabilité conditionnelle}
\ExplorerSection{Données}
Pour les exercices 1 et 2 : on choisit l’urne A avec probabilité \(\frac{1}{3},\) l’urne B avec probabilité \(\frac{2}{3}\). A contient 3 boules rouges et 1 bleue ; B contient 1 rouge et 3 bleues. On tire ensuite une boule au hasard dans l’urne choisie.\par
\vspace{5pt}
\ExplorerSection{Exercice 1}
Dans l’exemple, calculer P(A ∩ bleue).\par
\vspace{5pt}
\ExplorerSection{Exercice 2}
Calculer \(P(A|\mathrm{bleue})\).\par
\vspace{5pt}
\ExplorerSection{Exercice 3}
Dans une classe, 12 élèves sur 30 jouent aux échecs ; parmi eux 9 aiment les énigmes. Calculer P(énigmes|échecs). Peut-on en déduire P(échecs|énigmes) ?\par
\vspace{5pt}
\ExplorerSection{Exercice 4}
Deux événements A et R vérifient \(P(A)=0,4, P(R|A)=0,5\) et \(P(R|\mathrm{non} A)=0,5\). Calculer \(P(A|R)\).\par
\vspace{5pt}
\end{document}
